\section{Optimasi Portofolio dalam Ruang Biner $\mathbb{B}^N$}

\subsection{Formulasi Lagrangian Markowitz}
Dalam konteks pemilihan portofolio diskret, kita mendefinisikan semesta aset $N$ di mana keputusan investasi direpresentasikan oleh vektor profil strategi biner $\mathbf{x} \in \{0, 1\}^N$. Setiap elemen $x_i = 1$ menunjukkan bahwa aset $i$ disertakan dalam portofolio, sementara $x_i = 0$ menunjukkan eksklusi aset tersebut. Masalah optimasi Markowitz dalam ruang ini bertujuan untuk meminimalkan risiko varians portofolio $\mathbf{x}^T \Sigma \mathbf{x}$ sekaligus memaksimalkan ekspektasi imbal hasil $\bm{\mu}^T \mathbf{x}$ secara simultan. Fungsi objektif Lagrangian murni yang merepresentasikan dilema risiko-imbal hasil ini didefinisikan sebagai berikut:
\begin{equation}
    \min_{\mathbf{x} \in \mathbb{B}^N} \mathcal{L}(\mathbf{x}) = \sum_{i=1}^N \sum_{j=1}^N \sigma_{ij} x_i x_j - \lambda \sum_{i=1}^N \mu_i x_i
\end{equation}
di mana $\sigma_{ij}$ merupakan elemen matriks kovariansi yang mengukur interaksi risiko antar-aset dan $\lambda$ adalah parameter pengali Lagrange yang mengatur tingkat toleransi terhadap risiko. Penentuan nilai pengali Lagrange ini sangat bergantung pada profil risiko investor yang bersangkutan dalam menghadapi ketidakpastian pasar modal.

\subsection{Penyelarasan ke dalam Struktur Permainan Potensial}
Transisi dari optimasi terpusat menuju pemodelan berbasis teori permainan memerlukan penyelarasan parameter antara fungsi objektif global dan utilitas individual pemain. Kita mendefinisikan fungsi potensial global $\Phi(\mathbf{x})$ sebagai negatif dari fungsi Lagrangian yang telah diskalakan oleh konstanta penghindaran risiko $\gamma/2$. Dengan menetapkan hubungan analitis antara parameter risiko $\lambda = 2/\gamma$, fungsi potensial tersebut dapat dinyatakan kembali sebagai selisih antara total imbal hasil dan total risiko sistemik:
\begin{equation}
    \Phi(\mathbf{x}) = -\frac{\gamma}{2} \mathcal{L}(\mathbf{x}) = \sum_{i=1}^N \mu_i x_i - \frac{\gamma}{2} \sum_{i=1}^N \sum_{j=1}^N \sigma_{ij} x_i x_j
\end{equation}
Transformasi ini memastikan bahwa arah optimasi selaras dengan maksimisasi utilitas oleh agen-agen ekonomi yang bertindak secara rasional dalam sistem portofolio tersebut. Penyesuaian ini sangat krusial agar ekuilibrium yang dicapai dalam permainan potensial mencerminkan solusi optimal dari masalah Markowitz tradisional yang telah didiskritisasi.

\subsection{Dekomposisi Utilitas dan Bukti \textit{Exact Potential Game}}
Untuk mengidentifikasi fungsi \textit{payoff} dari perspektif pemain tunggal ke-$i$, fungsi potensial global didekomposisi dengan memisahkan kontribusi variabel keputusan $x_i$ dari variabel pemain lainnya $\mathbf{x}_{-i}$. Melalui pemecahan komponen jumlahan ganda dan pemanfaatan sifat simetri matriks kovariansi ($\sigma_{ij} = \sigma_{ji}$), fungsi $\Phi(\mathbf{x})$ dapat dijabarkan berdasarkan interaksi lokal aset $i$ sebagai berikut:
\begin{equation}
    \Phi(\mathbf{x}) = \mu_i x_i - \frac{\gamma}{2} \left( \sigma_{ii} x_i^2 + 2 \sum_{j \neq i} \sigma_{ij} x_i x_j \right) + C_{-i}
\end{equation}
di mana $C_{-i}$ merepresentasikan sekumpulan suku konstanta yang tidak bergantung pada tindakan pemain $i$. Konstanta ini mencakup imbal hasil dan risiko dari seluruh aset dalam semesta selain aset $i$ yang tetap konstan terhadap perubahan $x_i$.

Penerapan sifat fundamental variabel biner $x_i^2 = x_i$ memungkinkan kita untuk melakukan faktorisasi terhadap variabel keputusan $x_i$ dalam persamaan di atas. Dengan mengekstraksi $x_i$ dari seluruh suku yang relevan, kita memperoleh relasi eksplisit antara fungsi potensial dan utilitas individual $u_i(x_i, \mathbf{x}_{-i})$ sebagai berikut:
\begin{equation}
    \Phi(\mathbf{x}) = x_i \left( \mu_i - \frac{\gamma}{2} \sigma_{ii} - \gamma \sum_{j \neq i} \sigma_{ij} x_j \right) + C_{-i}
\end{equation}
Integrasi ini secara matematis membuktikan bahwa sistem Markowitz biner merupakan sebuah \textit{Exact Potential Game}. Relasi $\Phi(\mathbf{x}) = u_i(x_i, \mathbf{x}_{-i}) + C_{-i}$ menunjukkan bahwa setiap peningkatan utilitas individu yang dilakukan oleh satu pemain akan mengakibatkan peningkatan yang setara pada fungsi potensial global sistem.
